% Generated draft: requires mathematical and source review.
\documentclass[11pt]{article}
\usepackage[margin=25mm]{geometry}
\usepackage{fontspec}
\setmainfont{Latin Modern Roman}
\usepackage{amsmath,amssymb,mathtools,unicode-math}
\setmathfont{Latin Modern Math}
\usepackage{enumitem,needspace,xurl,hyperref}
\setlength{\parindent}{0pt}
\setlength{\parskip}{5pt}
\setlength{\emergencystretch}{3em}
\newcommand{\N}{\mathbb{N}}
\newcommand{\Z}{\mathbb{Z}}
\newcommand{\Q}{\mathbb{Q}}
\newcommand{\R}{\mathbb{R}}
\newcommand{\C}{\mathbb{C}}
\newcommand{\F}{\mathbb{F}}
\newcommand{\A}{\mathbb{A}}
\newcommand{\PP}{\mathbb{P}}
\newcommand{\E}{\mathbb{E}}
\newcommand{\eps}{\varepsilon}
\newcommand{\dd}{\,\mathrm{d}}
\begin{document}
\section*{UnsolvedMath ID: 30005360}
\textbf{Source problem code:} catalogue entry
\subsection{Definitions and mathematical statement}
Let \(X\) denote the set referred to in the supplied source, and let \((S_d)_{d\in D}\) denote its stated canonical partition into definable sets. Here \(D\) is only a placeholder for the unspecified index set. The asserted partition relation is
\[
X=\bigsqcup_{d\in D}S_d.
\]
The source record does not specify \(X\), its ambient structure, the meaning of definability, the index set, or the construction that makes the partition canonical.

The source asks: \emph{Do the canonical definable pieces \(S_d\) form a stratification of \(X\) in some classical sense?}

This question remains incompletely specified. No topology, differentiable or algebraic structure, or definition of stratification is supplied. In particular, the record does not identify any required regularity, local finiteness, frontier condition, or compatibility condition between pieces. The intended quantification over \(X\) and its ambient structures is also absent. Consequently, a precise stratification assertion cannot be extracted without additional source context; the displayed partition relation records the supplied premise and does not establish any stratification property.
\subsection{Short English statement}
The question asks whether a particular canonical partition into definable sets is a classical stratification. The supplied record omits the construction of the pieces, the ambient setting, and the intended stratification conditions, so it does not determine a precise conjecture.
\subsection{Sources}
\begin{itemize}
\item Supplied local UnsolvedMath registry record, permanent ID \(30005360\), title \emph{Stratification by Canonical Definable Partitions}, source problem code \texttt{catalogue entry}, in the Oberwolfach Reports Open Problems collection.
\item \emph{Model Theory: Combinatorics, Groups, Valued Fields and Neostability} (2023), Oberwolfach Reports, pp.\ 91--149, DOI \texttt{10.4171/owr/2023/2}. Supplied source and literature link: \url{https://doi.org/10.4171/owr/2023/2}. The supplied excerpt carries the report page marker 117. No individual proposer is identified in the record.
\item Literature assessment dated exactly \texttt{2026-08-22}: the supplied curation labels the question open because the required classical sense of stratification is unspecified and reports that no theorem resolving that exact formulation was identified. This is supplied dated curation, not an independently verified literature assessment. The background describes it as literature triage derived from the supplied Oberwolfach Reports curation and requests independent verification of time-sensitive status and cited claims.
\item The registry labels the statement \texttt{corrected\_verified}; its supplied verification note says that only formula delimiters, notation, grammar, and typography were repaired, with mathematical content and scope preserved. This label does not supply the missing mathematical context or constitute independent mathematical verification here.
\end{itemize}
\end{document}
